\section{Introducción y Antecedentes}
\label{sec:intro}

\subsection{Introducción}
En las últimas décadas, la física de la materia condensada ha experimentado una revolución con el descubrimiento de las fases topológicas de la materia, donde las propiedades globales de la estructura electrónica están protegidas contra perturbaciones locales y desorden por la simetría de inversión temporal u otras simetrías cristalinas~\cite{bansil2016colloquium}. 

Paralelamente, el estudio de heteroestructuras de baja dimensionalidad ha emergido como una de las plataformas más prometedoras para la ingeniería de materiales. Al combinar materiales con propiedades físicas distintas en una interfaz atómicamente plana, es posible inducir nuevos estados electrónicos y magnéticos que no están presentes en los materiales individuales. Este protocolo propone explorar la combinación de aleaciones de Heusler (conocidas por su alta polarización de espín y propiedades magnéticas controlables) y MXenes (carburos y nitruros de metales de transición bidimensionales con excelente conductividad y versatilidad superficial).

\subsection{Antecedentes}

\subsubsection{Aleaciones de Heusler}
Las aleaciones de Heusler son compuestos intermetálicos con estructura cúbica ($L2_1$ para Heusler completos o $C1_b$ para semi-Heusler) que exhiben una amplia variedad de propiedades físicas, desde ferromagnetismo de espín completo (half-metals) hasta superconductividad y comportamiento de aislantes topológicos~\cite{graf2011simple}. La posibilidad de sintonizar su brecha energética y la posición de sus niveles de energía mediante sustituciones químicas las convierte en candidatas ideales para acoplarse con materiales bidimensionales.

\subsubsection{MXenes}
Los MXenes son una familia de materiales bidimensionales descubiertos en 2011~\cite{naguib2011two}, con fórmula general $M_{n+1}X_n T_x$, donde $M$ representa un metal de transición temprana (como Ti, V, Cr), $X$ es carbono y/o nitrógeno, y $T_x$ representa grupos funcionales superficiales (como \ce{-O}, \ce{-F}, \ce{-OH}). Su gran área superficial, alta conductividad metálica e hidrofilia permiten una fácil integración en dispositivos nanoelectrónicos. Además, la presencia de metales de transición pesados en su estructura confiere un fuerte acoplamiento espín-órbita (SOC), el cual es un ingrediente clave para la inducción de estados topológicos no triviales.

\subsubsection{Heteroestructuras Heusler/MXene}
La interfase formada por una aleación de Heusler y un MXene ofrece un escenario sumamente rico para la modulación topológica. El fuerte acoplamiento de espín-órbita proveniente del MXene, combinado con el magnetismo y la polarización de espín de la aleación de Heusler, puede dar lugar a efectos físicos exóticos como el efecto Hall anómalo cuántico, la magnetorresistencia gigante o la generación de estados de borde topológicamente protegidos sintonizables. El control de la química y la estructura de la interfaz (ingeniería de interfaces) es, por lo tanto, el factor crucial para modular la topología de bandas en estas heteroestructuras.
